\chapter{Vertex covering in graphs with overlapping triangles}
\label{ch:covertriangles}

Chapters~\ref{ch:overlap} to~\ref{ch:metacomplex} measured one model on three
classes of network. The model was the Ising model, chosen because its couplings
are soft and its partition function is a sum rather than a search, so that what
the messages lose could be attributed to the structure and to nothing else.

This chapter asks whether the meta-complex repair survives a harder model. The
minimum vertex cover of Chapter~\ref{ch:cover} is a constraint problem: an
assignment is admissible or it is not, and a message carrying a forbidden value
carries a zero rather than a small number. Hard constraints are where message
passing is usually more fragile, not less, so the repair earning its exactness
here would say more than it did on a pairwise model.

\section{The model on a chygraph}
\label{sec:covermodel}

Put $x_{v}=1$ when atom $v$ is in the cover. Every bond must be covered, so a
complex carries the indicator that no bond lying inside it has both ends free,
\begin{equation}
  f_{A}(x_{A})=\prod_{\{i,j\}\subset A,\ \{i,j\}\in E}
    \bigl[1-(1-x_{i})(1-x_{j})\bigr],
  \label{eq:covercomplex}
\end{equation}
and a chemical potential $\mu$ is attached to each atom, $f_{v}(x_{v})
=e^{-\mu x_{v}}$, so that a large $\mu$ drives the cover to its minimum.

Two readings of Eq.~\eqref{eq:covercomplex} matter. On a clique it is
Sec.~\ref{sec:twoconstraints}'s statement that at most one member may be left
out, since any two members are adjacent. On a \emph{merged} complex it is not:
a meta-complex is a union of cliques rather than a clique, so only the bonds
actually present may be demanded, and the restriction to $\{i,j\}\in E$ is doing
work. Demanding all pairs would forbid covers that exist.

The recursion is the one of Chapters~\ref{ch:overlap} and~\ref{ch:metacomplex},
unchanged: messages between atoms and complexes, each complex summing over its
own interior, $\ln Z$ read off the Bethe free energy. Nothing about it is
specific to the Ising model --- Eq.~\eqref{eq:emitted} asks only that a
complex's interior can be enumerated --- and the same routine runs both, which
is what makes the two sets of numbers comparable.

\begin{figure}[t]
\centering
\begin{adjustbox}{max width=\linewidth}
\begin{tikzpicture}
% ---------------- left: leaf removal on the graph
\node[lb,anchor=south,align=center] at (0,2.05) {\emph{on the graph}};
\node[nd] (gu) at (-0.8,1.15) {};
\node[nd] (gw) at (-0.8,-0.15) {};
\node[nd] (gv) at (0.35,0.5) {};
\node[nd] (gp) at (-2.0,1.6) {};
\node[nd] (gq) at (-2.0,-0.6) {};
\node[lb,anchor=east] at (-0.95,1.15) {$u$};
\node[lb,anchor=east] at (-0.95,-0.15) {$w$};
\node[lb,anchor=west] at (0.5,0.5) {$v$};
\draw[lnk] (gu)--(gw) (gu)--(gv) (gw)--(gv) (gu)--(gp) (gw)--(gq);
\node[lb,anchor=north,align=center,text width=4.6cm] at (-0.6,-1.15)
  {$v$ has degree \emph{two}. No vertex here is a leaf, so the rule
   never fires and the whole triangle is core.};

\draw[lnk,dashed,black!35] (2.05,-1.6) -- (2.05,2.1);

% ---------------- right: the same object as a chygraph
\node[lb,anchor=south,align=center] at (5.1,2.05) {\emph{on the chygraph}};
\node[nd] (cu) at (4.3,1.15) {};
\node[nd] (cw) at (4.3,-0.15) {};
\node[nd] (cv) at (5.45,0.5) {};
\node[nd] (cp) at (3.1,1.6) {};
\node[nd] (cq) at (3.1,-0.6) {};
\node[lb,anchor=east] at (4.15,1.15) {$u$};
\node[lb,anchor=east] at (4.15,-0.15) {$w$};
\node[lb,anchor=west] at (5.6,0.5) {$v$};
\draw[lnk] (cu)--(cw) (cu)--(cv) (cw)--(cv) (cu)--(cp) (cw)--(cq);
\begin{scope}[on background layer]
  \node[cxb,fit=(cu)(cw)(cv)] {};
  \node[cx,fit=(cu)(cp)] {};
  \node[cx,fit=(cw)(cq)] {};
\end{scope}
\node[lb,anchor=north,align=center,text width=5.4cm] at (4.5,-1.15)
  {$v$ lies in \emph{one} complex, so it is a leaf one level up:
   cover $u$ and $w$, delete the complex, continue.};
\end{tikzpicture}
\end{adjustbox}
\caption{\textbf{Two leaf removals on one object.} Left, the rule of
Chapter~\ref{ch:cover}: find a vertex of degree one and take its neighbour.
Vertex $v$ has degree two, and neither $u$ nor $w$ is a leaf either, so the
algorithm stops at once and reports the triangle as core. Right, the same five
vertices as a chygraph, the triangle promoted to a complex. Now the test is
chy-degree rather than degree: $v$ belongs to one complex and to nothing else,
so it is needed nowhere but there. Since any cover of a clique leaves out at
most one member, and $v$ is the member that is free to be left out, $u$ and $w$
go into the cover and the complex is deleted. Promoting the triangle turned a
degree-two vertex into a leaf, which is the same trade the rest of the book
makes --- a loop becomes an interior --- appearing here as an algorithm that
finishes rather than one that stalls.}
\label{fig:twoleaf}
\end{figure}

\section{Where ordinary leaf removal stops}
\label{sec:ordinaryleaf}

Leaf removal needs a vertex of degree one. On the instances of Part~IV there
are almost none, and the algorithm stops essentially at the start.

Run on the graph, it consumes $3$ of the $120$: none of the thirty hyperbolic
instances, one of the thirty placed ones and two of the sixty real
neighbourhoods. On the other $117$ it halts on a core and returns nothing ---
not a poor cover, but no cover at all, since the guarantee is all-or-nothing.
The medians say why: the core it stops on is the whole graph on the hyperbolic
and real instances and $0.85$ of it on the placed ones.

That is not a defect of the algorithm. These are clustered graphs, and a
clustered graph has few leaves by construction: a vertex sitting in a triangle
has degree at least two, and the triangle is exactly what
Chapter~\ref{ch:chygraphs} promoted a complex to absorb. Chapter~\ref{ch:cover}
already said the core is where covering goes back to being hard, and on these
instances the core is everything.

\section{Leaf removal over complexes}
\label{sec:chyleaf}

The chygraph changes the test and nothing else. Where the graph rule asks for a
vertex of degree one, the chygraph rule asks for an atom of \emph{chy-degree}
one --- an atom belonging to a single complex (Fig.~\ref{fig:twoleaf}).

The move it licenses is safe, and the argument is Chapter~\ref{ch:cover}'s
exchange argument one level up. Let $v$ lie in the clique $A$ and in no other
complex. Any cover leaves out at most one member of $A$, since any two members
are adjacent. If a cover leaves out some $w\neq v$ then it contains $v$;
exchanging the two --- drop $v$, add $w$ --- is a cover of the same size, since
every neighbour of $v$ lies in $A$ and is now taken, and $w$ was taken anyway.
So some minimum cover leaves out $v$, and $A\setminus\{v\}$ may be put in the
cover and the complex deleted.

At $c=2$ this is the graph rule exactly: a complex of two atoms holding an atom
of chy-degree one is an edge hanging off a leaf. Above $c=2$ it is strictly
more. A vertex of degree $c-1$ inside a single clique is not a leaf on the
graph and is a leaf on the chygraph, so promoting the triangle converts the
degree-two vertices of Fig.~\ref{fig:twoleaf} into vertices the algorithm can
act on.

\paragraph{What that buys.} Table~\ref{tab:leafclasses} runs both rules on the
same $120$ instances.

\begin{table}[t]
\centering\small
\begin{tabular}{lrrrr}
\hline\hline
class & instances & ordinary & clique chygraph & merged\\
\hline
hyperbolic & $30$ & $0$ & $23$ & $27$\\
Karrer--Newman & $30$ & $1$ & $7$ & $19$\\
real & $60$ & $2$ & $40$ & $60$\\
\hline
total & $120$ & $3$ & $70$ & $106$\\
\hline\hline
\end{tabular}
\caption{\textbf{Instances solved exactly.} Number on which the algorithm
consumes the object and returns a provably minimum cover; on the rest it halts
on a core and returns nothing. Leaf removal on the graph manages $3$ of $120$.
The same rule applied to the clique chygraph manages $70$, and every one of the
three it already solved --- it solves $67$ the graph rule cannot and loses none.
Merging first raises it to $106$, for the reason Sec.~\ref{sec:mergedleaf}
gives. Verified against enumeration on every instance: where the algorithm
returns a cover, that cover is checked to be a cover and to have the minimum
size. Generated by \code{probe/cover\_leafremoval.py}.}
\label{tab:leafclasses}
\end{table}

Three of $120$ becomes seventy, and the improvement is one-directional: $67$
instances are solved by the chygraph rule and not by the graph rule, and none
the other way, so on this sample the chygraph rule strictly dominates. The gain
is largest where the clustering is highest, which is the point --- a triangle is
an obstruction to the graph algorithm and an interior to the chygraph one.

It is worth saying what this is not. The cover returned is not an approximation
that happens to be good. It is a minimum cover, certified by the same
all-or-nothing argument as Chapter~\ref{ch:cover}'s, and checked here against
enumeration on all $120$ instances. Where the chygraph rule finishes it has
solved an NP-hard problem exactly, on a clustered graph where the graph rule
could not start.

\section{Merging, and what it adds}
\label{sec:mergedleaf}

Chapter~\ref{ch:metacomplex}'s closure can be applied first, and
Table~\ref{tab:leafclasses}'s last column is the result: $106$ of $120$, and
all sixty real neighbourhoods.

The extra is not more leaf removal. A merged meta-complex is a union of cliques
and not a clique, so the exchange argument does not apply to it --- leaving out
one member of a meta-complex is not always allowed, because two of its members
need not be adjacent, and covering all but one of them can overcount. What the
closure supplies instead is \emph{isolated} complexes: a complex all of whose
atoms have chy-degree one shares nothing with the rest of the chygraph, so its
interior is enumerated and solved exactly.

Counting the moves separates the two. On the clique chygraph the exchange rule
fires $107$, $91$ and $97$ times across the three classes against $23$, $9$ and
$25$ enumerations. On the merged chygraph the exchange rule fires $20$, $46$ and
$20$ times and the enumerations rise to $30$, $22$ and $64$. So merging does not
extend the algorithm; it converts what the algorithm could not finish into
something small enough to enumerate --- which is the conclusion of
Sec.~\ref{sec:mergeerror} reached along a different route, and it carries the
same cost, $2^{c}$ on the largest meta-complex.

\section{Conclusions}
\label{sec:cover-tri-conclusions}

\paragraph{The chygraph earns its keep here more plainly than anywhere else in
the part.} Chapters~\ref{ch:overlap} to~\ref{ch:metacomplex} measured a
quantity the cavity method got wrong by a median of a few in $\ln Z$ and a
repair that removed it. This chapter measures an algorithm that either answers
or does not, and promoting the triangles takes it from answering $3$ of $120$
instances to answering $70$ --- exactly, with a certificate, on an NP-hard
problem.

\paragraph{The mechanism is the book's central move, in its simplest form.} A
vertex in a triangle has degree two and is not a leaf. The same vertex in a
promoted triangle has chy-degree one and is a leaf. Nothing was approximated to
achieve that; the exchange argument of Chapter~\ref{ch:cover} carries over to a
clique unchanged, and only the notion of degree was replaced.

\paragraph{It also shows what merging is for, and is not.} The closure raises
the count to $106$, but by supplying complexes small enough and isolated enough
to enumerate rather than by extending the rule, which cannot generalise to a
meta-complex --- a union of cliques is not a clique, and the exchange argument
needs a clique. That matches Sec.~\ref{sec:mergeerror}: merging buys exactness
by paying $2^{c}$, here as there.

\paragraph{And the hard-field cavity is the wrong tool for this job.}
Section~\ref{sec:hardfail} shows Eq.~\eqref{eq:vcz} returning density $1/2$ on
disjoint triangles where the truth is $2/3$, because it discards an $O(1)$ part
of the field that a complex makes essential. Leaf removal over complexes returns
$2/3$ on the same object, and does so with no field to discard. On the model
where an exact combinatorial rule exists, the combinatorial rule is better than
the message-passing one --- which is a caution about Part~IV's other chapters
rather than a result about this one.


\section{Checks}
\label{sec:covertri-checks}

\checkhead{Known results}
Leaf removal and its exchange argument are \citeauthor{bauer2001}'s, and
Chapter~\ref{ch:cover} derives them; the graph rule used for comparison is the
same routine Chapter~\ref{ch:cover} uses. The generalisation to a clique is
proved above and checked against enumeration: on every one of the $120$
instances where the rule terminates, the cover it returns is verified to be a
cover and to have exactly the minimum size, computed by enumeration over all
$2^{n}$ assignments.

\checkhead{New results}
The chygraph leaf-removal rule and its measurement are new here. Its anchors are
the two objects whose answers are known independently: a tree, consumed with the
exact cover; and disjoint triangles, where it returns $2/3$ against the
hard-field map's $1/2$. A bare cycle is left as a core, as it must be. The
comparison in Table~\ref{tab:leafclasses} is paired --- the same instance is
given to both rules --- and the domination is checked rather than assumed, by
counting the instances each solves that the other does not.

\checkhead{Partial results}
Only two moves are safe and only two are made. The exchange argument holds for a
clique and not for a merged meta-complex, so no move is attempted on a
non-clique complex still attached to the rest; an earlier version of this
calculation did attempt one, returned covers of the wrong size, and was caught
by the assertion against enumeration. What that leaves is a gap rather than a
result: whether some safe local move exists for a general complex is not
settled here, only that covering all but one atom is not it. The instances are
capped at twenty vertices because the verification enumerates, not because the
rule needs it --- leaf removal itself is linear, and the cap is the checking.


